\input includes/\format

\pdfinfo{
   /Title (Xuan Xuan Qi Jing)
   /Creator (TeX)
   /Author (Vit Brunner)
   /CreationDate (20261111233500)
}

~
\title
\parindent=0em
\baselineskip=1em

\vskip 4cm
\centerline{\cjktitle 玄玄棋經}
\bigskip
\centerline{Xuanxuan Qijing}
\bigskip
\centerline{(Gengen Gokyo)}
\vfill\break

\input includes/\format-motto

\parindent=1em
\bigskip
\noindent{\header Preface}
\medskip
\noindent {\it Gateway to All Marvels.} Xuanxuan Qijing is a classical Chinese collection compiled by Yan Defu and Yan Tianzhang around year 1347. Its perhaps more popular Japanese name is Gengen Gokyo.

This version is one of the more complete ones, containing 347 dan-level problems. All positions are shown as {\bf black to play}.

\bigskip

\bigskip

\vbox{
\noindent{\header Third edition note}
\medskip
\noindent Originally published in 2006, it's been twenty years. This booklet has been downloaded over a hundred thousand times. I never imagined the reach it would have.

\rightline{\it\me}
\rightline{\it November 2026}
}
\vfill\break

\input includes/\format-gotc
\input generated/\format-xxqj

\bye
